\begin{abstract}
This article defines a sieve sequence stage as the increasing sequence of
integers accepted by a finite prefix of the prime filters. The accepted
pattern is periodic modulo the product of those filters, so one finite
list of positive gaps reconstructs the entire infinite stage. We formally
verify the stage's strict increase, completeness, block-period shift,
gap-cycle reconstruction, and exact transition count. If the current head
is $h$, the current period has $T$ accepted values, and the current
modulus is $M$, then the expanded window of length $hM$ contains $hT$ old
survivors; exactly $T$ are multiples of $h$, leaving $T(h-1)$ survivors.
We also verify the local copy-or-merge rule for the next gaps. For each
real linear 2-gap, the two endpoint strikes occur at distinct lift
offsets, so exactly two of its $h$ lifts are destroyed and exactly $h-2$
keep both endpoints. Under explicit structural and period preconditions,
the resulting gap prefix agrees with the next linear specification.

The formal result has two explicit boundaries. First, next-head
primality is conditional on the square bound supplied mathematically by
Bertrand's postulate. Second, the article proves the semantic transition
between stages, while the direct construction of the next cycle from a
repeated and filtered current cycle is a separate open composition
problem. Accordingly, the article establishes a formally verified
finite-stage sieve specification and transition semantics, not a new
prime-sieving algorithm or a theorem about the persistence of any
particular prime gap in a prescribed short window.
\end{abstract}
